% Appendix A

\chapter{Questionário e Entrevistas} % Main appendix title

\label{AppendixA} % For referencing this appendix elsewhere, use \ref{AppendixA}

O conteúdo deste apêndice refere-se ao questionário em formato de entrevista realizado a agentes de investigação criminal da \gls{GNR} e da \gls{PSP} para validação do problema identificado.

\section{Síntese de Resultados}

A recolha de dados teve um carácter exploratório, focando-se na identificação de problemas práticos no acesso e análise de prova adquirida a parir de sistemas \gls{CCTV} durante a investigação criminal. As respostas obtidas junto dos \glspl{OPC} validam os desafios descritos na \autoref{sec:problema}, destacando-se as seguintes conclusões gerais:

\begin{itemize}
    \item \textbf{Conectividade e Acesso:} Confirma-se a utilização frequente de câmaras com cartão SIM/4G em redes não controladas e a dependência de dispositivos fora da rede institucional para o acesso inseuro às mesmas.
    \item \textbf{Fragmentação:} A generalidade dos inquiridos reporta a necessidade de utilizar múltiplas plataformas ou aplicações incompatíveis para aceder a diferentes equipamentos.
    \item \textbf{Ineficiência Operacional:} A análise de vídeo é descrita como um processo moroso (frequentemente demorando largos períodos de tempo para localizar e exportar eventos relevantes), assentando quase exclusivamente na revisão manual de gravações e \textit{playback} integral.
    \item \textbf{Restrições Institucionais:} Existe um forte consenso sobre a presença de restrições internas rigorosas quanto à instalação de novo software ou abertura de acessos diretos nas redes dos \glspl{OPC}.
    \item \textbf{Recetividade a Soluções Unificadas:} Os inquiridos demonstraram uma elevada recetividade a uma aplicação única que centralize o acesso às câmaras, integrando alertas e mecanismos de deteção automática.
\end{itemize}

\section{Questionário (ficheiro de respostas)}
\label{AppendixA:questionario_xlsx}

Para consulta e análise detalhada, o ficheiro Excel com as respostas do questionário encontra-se disponível numa versão estruturada e apresentável:

\begin{itemize}
    \item \textbf{Ficheiro (XLSX):} Disponível online em: \url{https://github.com/gustavocaiano/fuse/blob/main/docs/questionary/questionary_presentable.xlsx}
\end{itemize}
